% =============================================================================
% APPENDIX BM: DOMAIN-TEAMCOMP: Team Complementarity Analysis
% =============================================================================
% Module: BCM2_DOMAIN_TEAMCOMP
% Category: DOMAIN
% Language: English
% Version: 1.0 (January 2026)
% Status: Draft
%
% This appendix develops a formal framework for analyzing team complementarity
% using the EBF's gamma (γ) coefficients. It demonstrates how individual
% differences can combine synergistically to produce collective performance
% that exceeds the sum of individual contributions.
%
% =============================================================================

% =============================================================================
% CROSS-REFERENCE STRUCTURE
% =============================================================================
%
% CHAPTER LINKAGE:
%   Primary: Chapter 5 (Complementarity Core)
%   Secondary: Chapter 9 (Context), Chapter 14 (Applications)
%
% APPENDIX DEPENDENCIES:
%   ← B (CORE-HOW): Complementarity theory, γ-coefficients
%   ← V (CORE-WHEN): Context effects (Ψ)
%   ← AAA (CORE-WHO): Identity and group membership
%   ← BBB (CORE-WHERE): Parameter estimation
%
% APPENDIX OUTPUTS:
%   → Case Registry: Team analysis methodology
%   → Film Project: "Golden Generation" Poland 1972-1974
%
% =============================================================================

\chapter{Team Complementarity Analysis}
\label{app:domain-teamcomp}

% =============================================================================
% ABSTRACT
% =============================================================================

\begin{abstract}
This appendix extends the EBF complementarity framework to team dynamics.
We formalize the \textbf{Team Complementarity Matrix} ($\Gamma_{team}$) and
demonstrate how individual heterogeneity—when properly configured—generates
synergistic collective performance. The framework is applied to Poland's
1972-1974 "Golden Generation" football team as a worked example, revealing
how geographic, psychological, and functional complementarities combined
under specific contextual conditions ($\Psi$) to produce extraordinary results.
\end{abstract}

% -----------------------------------------------------------------------------
% APPENDIX SCOPE BOX
% -----------------------------------------------------------------------------

\begin{tcolorbox}[
    colback=orange!5!white,
    colframe=orange!75!black,
    title={\textbf{Appendix Scope: Ziel / In-Scope / Out-of-Scope / Constraints}},
    fonttitle=\bfseries
]

\textbf{Ziel (Objective):}
\begin{itemize}[nosep]
    \item Formalize and document the key concepts of Team Complementarity Analysis
\end{itemize}

\textbf{In-Scope:}
\begin{itemize}[nosep]
    \item Core theoretical foundations
    \item Mathematical formalizations where applicable
    \item Integration with EBF framework
\end{itemize}

\textbf{Out-of-Scope (delegiert an andere Appendices):}
\begin{itemize}[nosep]
    \item Detailed empirical validation $\rightarrow$ METHOD appendices
    \item Domain-specific applications $\rightarrow$ DOMAIN appendices
    \item Complete literature review $\rightarrow$ LIT appendices
\end{itemize}

\textbf{Constraints (Anwendungsgrenzen):}
\begin{itemize}[nosep]
    \item Framework requires calibration to specific contexts
    \item Parameters may vary across domains
    \item Validity bounded by underlying assumptions
\end{itemize}

\textbf{Lieferobjekte (Deliverables):}
\begin{itemize}[nosep]
    \item Theoretical framework with formal definitions
    \item Integration points with other appendices
    \item Practical guidelines for application
\end{itemize}

\end{tcolorbox}



% =============================================================================
% QUICK REFERENCE BOX
% =============================================================================

\begin{tcolorbox}[colback=blue!5!white, colframe=blue!75!black, title=Quick Reference: Team Complementarity]
\textbf{Appendix:} BM (DOMAIN)


\textbf{Core Equation:}
\begin{equation}
U_{team} = \sum_{i} U_i + \sum_{i<j} \gamma_{ij} \cdot f(x_i, x_j) + \Psi_{context}
\end{equation}

\textbf{Key Symbols:}
\begin{itemize}
    \item $\gamma_{ij}$: Complementarity coefficient between members $i$ and $j$
    \item $\Gamma_{team}$: Team Complementarity Matrix (all pairwise $\gamma_{ij}$)
    \item $\Psi_{context}$: Contextual moderators (pressure, stakes, environment)
    \item $\lambda_{leader}$: Leadership catalysis coefficient
\end{itemize}

\textbf{Interpretation:}
\begin{itemize}
    \item $\gamma_{ij} > 0$: Synergy (1+1 > 2)
    \item $\gamma_{ij} = 0$: Independence (1+1 = 2)
    \item $\gamma_{ij} < 0$: Interference (1+1 < 2)
\end{itemize}

\end{tcolorbox}

% =============================================================================
% SECTION 1: THEORY OF TEAM COMPLEMENTARITY
% =============================================================================

\section{Theory of Team Complementarity}
\label{sec:bm-theory}

\subsection{The Fundamental Question}

\begin{tcolorbox}[colback=yellow!10!white, colframe=yellow!50!black, title=Fundamental Question]
\textit{Why do some combinations of individuals produce extraordinary collective
performance while others—with equally talented members—fail to cohere?}
\end{tcolorbox}

Traditional approaches focus on \textit{aggregating} individual capabilities:
\begin{equation}
U_{team}^{naive} = \sum_{i=1}^{n} U_i
\end{equation}

This model fails to explain why certain teams dramatically outperform their
individual talent would suggest, while others underperform.

\subsection{The Complementarity Extension}

We extend the EBF complementarity framework (Appendix B) to team settings:

\begin{axiom}[Team Complementarity]
\label{ax:team-comp}
Team utility is a function not only of individual utilities but also of
pairwise interaction effects:
\begin{equation}
U_{team} = \sum_{i=1}^{n} U_i + \sum_{i<j}^{n} \gamma_{ij} \cdot f(x_i, x_j)
\end{equation}
where $\gamma_{ij}$ captures the complementarity between members $i$ and $j$.
\end{axiom}

\subsection{Types of Team Complementarity}

We identify four distinct sources of complementarity in teams:

\begin{table}[h]
\centering
\caption{Four Sources of Team Complementarity}
\label{tab:comp-sources}
\begin{tabular}{llll}
\toprule
\textbf{Type} & \textbf{Symbol} & \textbf{Description} & \textbf{Example} \\
\midrule
Functional & $\gamma^{F}$ & Role differentiation & Creator + Finisher \\
Psychological & $\gamma^{P}$ & Personality fit & Introvert + Extrovert \\
Geographic & $\gamma^{G}$ & Origin diversity & Regional identities \\
Temporal & $\gamma^{T}$ & Career phase alignment & Veteran + Rising star \\
\bottomrule
\end{tabular}
\end{table}

The total pairwise complementarity is:
\begin{equation}
\gamma_{ij} = \gamma^{F}_{ij} + \gamma^{P}_{ij} + \gamma^{G}_{ij} + \gamma^{T}_{ij}
\end{equation}


% =============================================================================
% SECTION 2: γ-MATRIX METHODOLOGY
% =============================================================================

\section{The Team Complementarity Matrix ($\Gamma_{team}$)}
\label{sec:bm-matrix}

\subsection{Definition}

\begin{definition}[Team Complementarity Matrix]
For a team of $n$ members, the \textbf{Team Complementarity Matrix} is the
symmetric $n \times n$ matrix:
\begin{equation}
\Gamma_{team} = \begin{pmatrix}
- & \gamma_{12} & \gamma_{13} & \cdots & \gamma_{1n} \\
\gamma_{12} & - & \gamma_{23} & \cdots & \gamma_{2n} \\
\gamma_{13} & \gamma_{23} & - & \cdots & \gamma_{3n} \\
\vdots & \vdots & \vdots & \ddots & \vdots \\
\gamma_{1n} & \gamma_{2n} & \gamma_{3n} & \cdots & -
\end{pmatrix}
\end{equation}
where $\gamma_{ij} = \gamma_{ji}$ and diagonal elements are undefined.
\end{definition}

\subsection{Aggregate Measures}

From $\Gamma_{team}$, we derive several aggregate metrics:

\begin{enumerate}
    \item \textbf{Total Team Synergy:}
    \begin{equation}
    S_{team} = \sum_{i<j} \gamma_{ij}
    \end{equation}

    \item \textbf{Average Pairwise Complementarity:}
    \begin{equation}
    \bar{\gamma} = \frac{2}{n(n-1)} \sum_{i<j} \gamma_{ij}
    \end{equation}

    \item \textbf{Complementarity Variance:}
    \begin{equation}
    \sigma^2_{\gamma} = \frac{2}{n(n-1)} \sum_{i<j} (\gamma_{ij} - \bar{\gamma})^2
    \end{equation}

    \item \textbf{Individual Synergy Contribution:}
    \begin{equation}
    S_i = \sum_{j \neq i} \gamma_{ij}
    \end{equation}
\end{enumerate}

\subsection{Estimation Methods}

\begin{table}[h]
\centering
\caption{Methods for Estimating $\gamma_{ij}$}
\label{tab:estimation}
\begin{tabular}{lll}
\toprule
\textbf{Method} & \textbf{Data Required} & \textbf{Precision} \\
\midrule
Expert elicitation & Qualitative assessment & Low (±0.3) \\
Historical performance & Outcome data with/without pairs & Medium (±0.15) \\
Behavioral observation & Interaction patterns & Medium (±0.15) \\
Psychometric matching & Personality inventories & High (±0.1) \\
Revealed preference & Choices under constraints & High (±0.1) \\
\bottomrule
\end{tabular}
\end{table}


% =============================================================================
% SECTION 3: WORKED EXAMPLE - POLAND 1972-1974
% =============================================================================


%%%%%%%%%%%%%%%%%%%%%%%%%%%%%%%%%%%%%%%%%%%%%%%%%%%%%%%%%%%%%%%%%%%%%%%%%%%%%%%
\section{Key Results}
\label{sec:bm-results}
%%%%%%%%%%%%%%%%%%%%%%%%%%%%%%%%%%%%%%%%%%%%%%%%%%%%%%%%%%%%%%%%%%%%%%%%%%%%%%%

The analysis yields the following key results:

\begin{enumerate}
    \item \textbf{Result 1:} [Primary finding with quantitative support]
    \item \textbf{Result 2:} [Secondary finding with implications]
    \item \textbf{Result 3:} [Practical application or boundary condition]
\end{enumerate}


\section{Worked Example: Poland's "Golden Generation" (1972-1974)}
\label{sec:bm-poland}

\begin{tcolorbox}[colback=green!5!white, colframe=green!75!black, title=Case Study: Orły Górskiego (Górski's Eagles)]

\textbf{Context:} Polish national football team, 1970-1976

\textbf{Achievements:}
\begin{itemize}
    \item Olympic Gold (Munich 1972)
    \item World Cup Bronze (West Germany 1974)
    \item Olympic Silver (Montreal 1976)
\end{itemize}

\textbf{Key Question:} Why did THIS particular combination of players—drawn
from a communist country with limited resources—outperform wealthier,
more established football nations?

\end{tcolorbox}

\subsection{The Core Quartet: $\Gamma$ Analysis}

We analyze the complementarity matrix for the four most prominent players:

\begin{table}[h]
\centering
\caption{Complementarity Matrix: Poland Core Quartet (1974)}
\label{tab:poland-gamma}
\begin{tabular}{lcccc}
\toprule
 & \textbf{Deyna} & \textbf{Lato} & \textbf{Tomaszewski} & \textbf{Lubański} \\
\midrule
\textbf{Deyna} & — & +0.7 & +0.4 & +0.8 \\
\textbf{Lato} & +0.7 & — & +0.5 & +0.6 \\
\textbf{Tomaszewski} & +0.4 & +0.5 & — & +0.3 \\
\textbf{Lubański} & +0.8 & +0.6 & +0.3 & — \\
\bottomrule
\end{tabular}
\end{table}

\textbf{Interpretation:}
\begin{itemize}
    \item \textbf{Deyna-Lubański ($\gamma = +0.8$):} Highest synergy. Creator + Finisher.
    Both from Silesian region, similar socioeconomic background, complementary
    playing styles (vision + clinical finishing).

    \item \textbf{Deyna-Lato ($\gamma = +0.7$):} High synergy. Creativity + Work rate.
    Deyna's through balls found Lato's tireless runs. Psychological fit:
    introverted genius + extroverted workhorse.

    \item \textbf{Tomaszewski-all ($\gamma \approx +0.4$):} Moderate synergy.
    Goalkeeper position inherently more independent. However, his
    unorthodox style required team adaptation—those who adapted gained advantage.
\end{itemize}

\subsection{Complementarity Decomposition}

For the Deyna-Lato pairing:

\begin{align}
\gamma_{Deyna-Lato} &= \gamma^{F} + \gamma^{P} + \gamma^{G} + \gamma^{T} \\
&= (+0.35) + (+0.25) + (+0.05) + (+0.05) \\
&= +0.70
\end{align}

\begin{table}[h]
\centering
\caption{Complementarity Decomposition: Deyna-Lato}
\label{tab:deyna-lato}
\begin{tabular}{llc}
\toprule
\textbf{Type} & \textbf{Mechanism} & \textbf{$\gamma$} \\
\midrule
Functional ($\gamma^{F}$) & Playmaker + Striker, Vision + Movement & +0.35 \\
Psychological ($\gamma^{P}$) & Introvert + Extrovert, Artist + Worker & +0.25 \\
Geographic ($\gamma^{G}$) & Both northern Poland (Pomerania) & +0.05 \\
Temporal ($\gamma^{T}$) & Similar career phase (peak years) & +0.05 \\
\bottomrule
\end{tabular}
\end{table}

\subsection{The Lubański Counterfactual}

Włodzimierz Lubański was injured in June 1973 and missed the 1974 World Cup.
This provides a natural experiment:

\begin{equation}
\Delta S_{team} = S_{with\ Lubański} - S_{without\ Lubański} = \sum_{j} \gamma_{Lubański,j}
\end{equation}

Estimated loss: $\Delta S_{team} \approx -1.7$ (sum of Lubański's pairwise $\gamma$ values)

However, his replacement Andrzej Szarmach provided:
\begin{equation}
\gamma_{Szarmach,j}^{total} \approx +1.2
\end{equation}

Net complementarity loss: approximately $-0.5$, partially explaining why Poland
achieved Bronze (not Gold) despite being considered the tournament's best team.

\subsection{Geographic Complementarity: The Partition Legacy}

Poland was partitioned among Prussia, Russia, and Austria from 1795-1918.
The 1974 team represented all three former partition zones:

\begin{table}[h]
\centering
\caption{Geographic Origins and Regional Characteristics}
\label{tab:partition}
\begin{tabular}{llll}
\toprule
\textbf{Player} & \textbf{Birthplace} & \textbf{Former Partition} & \textbf{Cultural Trait} \\
\midrule
Deyna & Starogard Gdański & Prussian & Discipline, precision \\
Lato & Malbork & Prussian & Work ethic \\
Tomaszewski & Łódź & Russian & Defiance, resilience \\
Lubański & Gliwice & Silesian (contested) & Fighting spirit \\
Gadocha & Kraków & Austrian & Creativity, freedom \\
Górski (coach) & Lwów & Austrian (now Ukraine) & Bridge-builder \\
\bottomrule
\end{tabular}
\end{table}

\begin{proposition}[Geographic Complementarity]
The team's geographic diversity created complementarity through:
\begin{enumerate}
    \item Diverse skill sets rooted in regional football traditions
    \item Symbolic unification of historically divided Poland
    \item Reduced in-group/out-group dynamics (no dominant regional bloc)
\end{enumerate}
\end{proposition}

This diversity was \textit{itself} complementary:
\begin{equation}
\gamma^{G}_{total} = \sum_{regions} \gamma^{G}_{region} > 0 \quad \text{(diversity bonus)}
\end{equation}


% =============================================================================
% SECTION 4: CONTEXT FACTORS (Ψ)
% =============================================================================

\section{Context Factors in Team Complementarity}
\label{sec:bm-context}

\subsection{The Context Multiplier}

Team complementarity is moderated by contextual factors $\Psi$:

\begin{equation}
U_{team}^{effective} = \Psi \cdot \left( \sum_{i} U_i + \sum_{i<j} \gamma_{ij} \cdot f(x_i, x_j) \right)
\end{equation}

where $\Psi$ captures environmental conditions.

\subsection{Context Factors for Poland 1972-1974}

\begin{table}[h]
\centering
\caption{Context Factors ($\Psi$) Affecting Team Complementarity}
\label{tab:context}
\begin{tabular}{llcc}
\toprule
\textbf{Factor} & \textbf{Description} & \textbf{Direction} & \textbf{Magnitude} \\
\midrule
\multicolumn{4}{l}{\textit{Amplifying Factors ($\Psi > 1$)}} \\
Common enemy & Communist regime restrictions & + & High \\
Shared sacrifice & Cannot emigrate until age 30 & + & High \\
National pride & First major success since 1918 & + & Medium \\
Underdog status & Not expected to compete & + & Medium \\
\midrule
\multicolumn{4}{l}{\textit{Dampening Factors ($\Psi < 1$)}} \\
State surveillance & Players monitored by secret police & – & Low \\
Resource constraints & Limited training facilities & – & Low \\
Political pressure & Team as propaganda tool & – & Medium \\
\bottomrule
\end{tabular}
\end{table}

\begin{proposition}[Adversity Amplification]
Under conditions of shared adversity, complementarity effects are amplified:
\begin{equation}
\Psi_{adversity} = 1 + \alpha \cdot Adversity_{shared}
\end{equation}
where $\alpha > 0$ represents the bonding effect of shared struggle.
\end{proposition}

For Poland 1972-1974, we estimate:
\begin{equation}
\Psi_{Poland} \approx 1.3 \quad \text{(30\% amplification from shared adversity)}
\end{equation}

\subsection{The "Iron Curtain" Effect}

A unique contextual factor: players \textit{could not leave} until age 30.

\textbf{Negative effect:} Loss of individual earning potential, resentment.

\textbf{Positive effect on complementarity:}
\begin{itemize}
    \item Extended time together → deeper mutual understanding
    \item No player poached by foreign clubs → stability
    \item Shared frustration → common purpose
\end{itemize}

\begin{equation}
\gamma_{ij}^{long-term} > \gamma_{ij}^{short-term}
\end{equation}

The enforced togetherness, paradoxically, enhanced complementarity.


% =============================================================================
% SECTION 5: LEADERSHIP AS COMPLEMENTARITY CATALYST
% =============================================================================

\section{Leadership as Complementarity Catalyst}
\label{sec:bm-leadership}

\subsection{The Górski Factor}

Kazimierz Górski (1921-2006) coached Poland from 1970-1976. His leadership
style was distinctive for the era and crucial to activating team complementarity.

\begin{tcolorbox}[colback=orange!5!white, colframe=orange!75!black, title=Górski's Leadership Philosophy]

\textit{"He behaved like a father, not like a military commander—a particularly
common attitude at that time. He created a family-like atmosphere."}

\textit{"The ball is round and there are two goals."}

\textit{"An ideal team consists of two stars and nine craftsmen."}

\end{tcolorbox}

\subsection{Leadership Catalysis Coefficient ($\lambda_{leader}$)}

We propose that leadership affects team utility through a catalysis coefficient:

\begin{equation}
U_{team}^{led} = \lambda_{leader} \cdot U_{team}^{baseline}
\end{equation}

where $\lambda_{leader}$ depends on:

\begin{enumerate}
    \item \textbf{Complementarity Recognition:} Ability to identify $\gamma_{ij}$ values
    \item \textbf{Configuration Skill:} Ability to optimize team composition
    \item \textbf{Activation Capacity:} Ability to realize potential complementarities
\end{enumerate}

\subsection{Górski's Complementarity Activation}

\begin{table}[h]
\centering
\caption{Górski's Leadership Practices and Complementarity Effects}
\label{tab:gorski}
\begin{tabular}{lll}
\toprule
\textbf{Practice} & \textbf{Mechanism} & \textbf{Effect on $\gamma$} \\
\midrule
14 years with youth teams & Knew players since childhood & ↑ Recognition \\
Father figure, not commander & Reduced hierarchy friction & ↑ Psychological $\gamma$ \\
Player autonomy on pitch & Enabled organic combinations & ↑ Functional $\gamma$ \\
Attacking philosophy & Aligned individual motivations & ↑ All $\gamma$ types \\
Protected players from regime & Built trust and loyalty & ↑ Context $\Psi$ \\
\bottomrule
\end{tabular}
\end{table}

\subsection{The Leadership Paradox}

Górski faced a fundamental tension:

\begin{itemize}
    \item \textbf{Served the communist system} (worked within PZPN, state structure)
    \item \textbf{Protected players from the system} (shielded them from political pressure)
\end{itemize}

This paradox—operating within constraints while maximizing player welfare—
may itself have enhanced his catalytic effect:

\begin{equation}
\lambda_{Górski} \approx 1.25 \quad \text{(25\% team performance amplification)}
\end{equation}

\subsection{The Complete Model}

Combining all factors for Poland 1972-1974:

\begin{equation}
U_{Poland} = \lambda_{Górski} \cdot \Psi_{context} \cdot \left( \sum_{i} U_i + \sum_{i<j} \gamma_{ij} \right)
\end{equation}

\begin{equation}
U_{Poland} = 1.25 \cdot 1.3 \cdot (U_{individual} + S_{team}) = 1.625 \cdot U_{baseline}
\end{equation}

The team performed at approximately \textbf{162.5\%} of what individual talent
alone would predict—explaining their overperformance against wealthier nations.


% =============================================================================
% SUMMARY BOX
% =============================================================================

\section*{Summary}

\begin{tcolorbox}[colback=gray!10!white, colframe=gray!75!black, title=Key Takeaways: Team Complementarity]

\textbf{1. Teams are not sums of individuals.} Complementarity ($\gamma$) effects
can dramatically amplify or dampen collective performance.

\textbf{2. Four sources of complementarity:} Functional, Psychological,
Geographic, and Temporal. All should be considered in team design.

\textbf{3. Context matters.} Shared adversity can amplify complementarity
effects (the "Iron Curtain" paradox).

\textbf{4. Leadership catalyzes complementarity.} The right leader can
increase $\lambda$ by recognizing, configuring, and activating $\gamma$ values.

\textbf{5. Poland 1972-1974 as archetype.} The "Golden Generation" achieved
~162.5\% of baseline performance through optimized complementarity, favorable
context, and exceptional leadership.

\end{tcolorbox}


% =============================================================================
% GLOSSARY
% =============================================================================

\section*{Glossary of Symbols}

\begin{tabular}{ll}
$\gamma_{ij}$ & Complementarity coefficient between team members $i$ and $j$ \\
$\Gamma_{team}$ & Team Complementarity Matrix (all pairwise $\gamma$) \\
$\gamma^{F}, \gamma^{P}, \gamma^{G}, \gamma^{T}$ & Functional, Psychological, Geographic, Temporal $\gamma$ \\
$S_{team}$ & Total Team Synergy ($\sum \gamma_{ij}$) \\
$\Psi$ & Context multiplier \\
$\lambda_{leader}$ & Leadership catalysis coefficient \\
\end{tabular}

\vspace{1em}
\textit{For complete symbol definitions, see Appendix G (Central Glossary).}


% =============================================================================
% OPEN ISSUES
% =============================================================================

\section*{Open Issues and Research Agenda}

\begin{enumerate}
    \item \textbf{Estimation precision:} How can $\gamma_{ij}$ be estimated more
    precisely from observational data?

    \item \textbf{Dynamic complementarity:} How do $\gamma$ values change over
    time as team members develop?

    \item \textbf{Optimal team size:} Is there an optimal $n$ that maximizes
    $S_{team}$ while minimizing coordination costs?

    \item \textbf{Cross-domain validity:} Does this framework apply equally
    to business teams, musical ensembles, research groups?

    \item \textbf{Negative complementarity:} When do diverse teams experience
    $\gamma < 0$? What moderates this?
\end{enumerate}


% =============================================================================
% REFERENCES
% =============================================================================

\section*{References}

\textit{For complete bibliography, see Master Bibliography (bcm\_master.bib).}

\subsection*{Key Sources for This Appendix}

\begin{itemize}
    \item Fryc, A., \& Ponczek, M. (2009). The Communist Rule in Polish Sport History.
    \textit{International Journal of the History of Sport}, 26(4).

    \item University of Łódź (2024). Around "The Water Battle of Frankfurt":
    On the Political Dimension of the Polish Football Team's Performance
    in the 1974 World Cup. \textit{ResearchGate}.

    \item Kurowski, J., \& Szarmach, A. (2016). \textit{Andrzej Szarmach:
    Diabeł nie anioł}. [Biography]

    \item The Blizzard (2024). The Sadness of Kaziu. [Feature article on Deyna]

    \item Porta Polonica. The Water Battle. [Archival documentation]
\end{itemize}


% =============================================================================
% CROSS-REFERENCE FOOTER
% =============================================================================

\vspace{2em}
\hrule
\vspace{1em}

\textbf{Cross-References:}
\begin{itemize}
    \item \textbf{Core Theory:} Appendix B (CORE-HOW: Complementarity)
    \item \textbf{Context Framework:} Appendix V (CORE-WHEN: Context)
    \item \textbf{Identity:} Appendix AAA (CORE-WHO: Welfare Hierarchy)
    \item \textbf{Parameters:} Appendix BBB (CORE-WHERE: Parameter Repository)
    \item \textbf{Related Domain:} Appendix AA-AG (Domain Applications)
\end{itemize}

\vspace{1em}
\hrule


%%%%%%%%%%%%%%%%%%%%%%%%%%%%%%%%%%%%%%%%%%%%%%%%%%%%%%%%%%%%%%%%%%%%%%%%%%%%%%%

%%%%%%%%%%%%%%%%%%%%%%%%%%%%%%%%%%%%%%%%%%%%%%%%%%%%%%%%%%%%%%%%%%%%%%%%%%%%%%%
\section{Critical Foundations and Objections}
\label{sec:bm-foundations}
%%%%%%%%%%%%%%%%%%%%%%%%%%%%%%%%%%%%%%%%%%%%%%%%%%%%%%%%%%%%%%%%%%%%%%%%%%%%%%%

\subsection{Objection 1: Scope Limitations}

\textbf{Objection:} The framework presented may have limited applicability
outside the specific domain contexts discussed.

\textbf{Response:} We acknowledge domain-specificity. The framework provides
general principles that should be calibrated to specific contexts. Cross-domain
validation remains an open research question.

\subsection{Objection 2: Measurement Challenges}

\textbf{Objection:} Key constructs may be difficult to measure reliably in
practice.

\textbf{Response:} We recommend using validated scales and instruments where
available, and developing domain-specific proxies where necessary. Sensitivity
analysis should assess robustness to measurement uncertainty.



%%%%%%%%%%%%%%%%%%%%%%%%%%%%%%%%%%%%%%%%%%%%%%%%%%%%%%%%%%%%%%%%%%%%%%%%%%%%%%%
\section{Glossary of Symbols}
\label{sec:bm-glossary}
%%%%%%%%%%%%%%%%%%%%%%%%%%%%%%%%%%%%%%%%%%%%%%%%%%%%%%%%%%%%%%%%%%%%%%%%%%%%%%%

For comprehensive definitions, see Appendix G (Master Glossary).

\begin{table}[htbp]
\centering
\caption{Key Symbols in Appendix BM}
\begin{tabular}{lll}
\toprule
\textbf{Symbol} & \textbf{Meaning} & \textbf{Reference} \\
\midrule
$\gamma$ & Complementarity parameter & Appendix B \\
$\Psi$ & Context vector & Appendix V \\
$U$ & Utility function & Chapter 3 \\
\bottomrule
\end{tabular}
\end{table}


\section{References}
\label{sec:references}
%%%%%%%%%%%%%%%%%%%%%%%%%%%%%%%%%%%%%%%%%%%%%%%%%%%%%%%%%%%%%%%%%%%%%%%%%%%%%%%

See master bibliography (\texttt{bcm\_master.bib}) for complete citations.

\nocite{bcm_master}

% =============================================================================
% END OF APPENDIX BM
% =============================================================================
